\documentclass[11pt,spanish,a4paper,]{article}
\usepackage[T1]{fontenc}
\usepackage{lmodern}
\usepackage{amssymb,amsmath}
\usepackage{ifxetex,ifluatex}
\usepackage{fixltx2e} % provides \textsubscript
\usepackage{graphicx}
\usepackage{wrapfig}
\usepackage{booktabs}
\usepackage[usenames,dvipsnames]{color}
\usepackage{lastpage}
%\usepackage{ccicons}
% use upquote if available, for straight quotes in verbatim environments
\IfFileExists{upquote.sty}{\usepackage{upquote}}{}
\ifnum 0\ifxetex 1\fi\ifluatex 1\fi=0 % if pdftex
  \usepackage[utf8]{inputenc}
\usepackage{titlesec}

\titleformat{\section}
{\color{red}\normalfont\Large\bfseries}
{\color{red}\thesection}{1em}{}
\newcommand{\doubleSignature}[4]{%
  \parbox{\textwidth}{
    #3, a #4\\
    \vspace{3cm}

    \parbox{7cm}{
      \centering
      \rule{6cm}{1pt}\\
       Fmdo: #1 
    }
    \hfill
    \parbox{7cm}{
      \centering
      \rule{6cm}{1pt}\\
      Fmdo: #2
    }
  }
}
\else % if luatex or xelatex
  \usepackage{fontspec}
  \ifxetex
    \usepackage{xltxtra,xunicode}
  \fi
  \defaultfontfeatures{Mapping=tex-text,Scale=MatchLowercase}
  \newcommand{\euro}{€}
    \setmainfont{Ubuntu}
\fi
% use microtype if available
\IfFileExists{microtype.sty}{\usepackage{microtype}}{}
\usepackage[margin=1in]{geometry}
\ifxetex
  \usepackage[setpagesize=false, % page size defined by xetex
              unicode=false, % unicode breaks when used with xetex
              xetex]{hyperref}
\else
  \usepackage[unicode=true]{hyperref}
\fi
\hypersetup{breaklinks=true,
            bookmarks=true,
            pdfauthor={},
            pdftitle={GuifiBaix: Contrato de mantenimiento y servicios, Alberta Gijón},
            colorlinks=true,
            urlcolor=blue,
            linkcolor=magenta,
            pdfborder={0 0 0}}
\urlstyle{same}  % don't use monospace font for urls
\setlength{\parindent}{0pt}
\setlength{\parskip}{6pt plus 2pt minus 1pt}
\setlength{\emergencystretch}{3em}  % prevent overfull lines
\setcounter{secnumdepth}{0}
\ifxetex
  \usepackage{polyglossia}
  \setmainlanguage{spanish}
\else
  \usepackage[spanish]{babel}
\fi
\renewcommand*{\familydefault}{\sfdefault}

\usepackage{xcolor}
\usepackage{fancyhdr}
\pagestyle{fancy}
\lhead{\includegraphics[height=10mm]{guifibaix-logo}}
\cfoot{\thepage / \pageref{LastPage}}
\rhead{Contrato de mantenimiento y servicios}
\title{\raisebox{-.4\height}{\includegraphics[width=4cm]{guifibaix-logo}} \hspace{1cm} \color{tuscanred}Contrato de mantenimiento y servicios}
\author{}
\date{}

\usepackage{sectsty}
\definecolor{tuscanred}{rgb}{0.51, 0.21, 0.21}
\allsectionsfont{\color{tuscanred}}  % sets colour of sections


\begin{document}

\maketitle

En Sant Joan Despí, a 20 de febrero de 2014

\subsection{REUNIDOS}\label{reunidos}

DE UNA PARTE, \textbf{Alberta Gijón} mayor de edad, con N.I.F. número
\textbf{12345678V}, en adelante, la ``CLIENTA'', con domicilio en
\textbf{Sant Joan Despí}, \textbf{C/Rue Percebe, 13, 4o 3a}, C.P.
\textbf{08970}.

DE OTRA PARTE, \textbf{David García Garzón}, con N.I.F. número
\textbf{36517097C} y en nombre y representación de \textbf{GUIFIBAIX
SCCL}, en adelante, la ``PROVEEDORA'', con domicilio en \textbf{Sant
Joan Despí}, \textbf{C/Major 22, 3}, C.P. \textbf{08970} y C.I.F.
\textbf{F66576380}.

La CLIENTA y la PROVEEDORA, en adelante, podrán ser denominadas,
individualmente, ``la Parte'' y, conjuntamente, ``las Partes'',
reconociéndose mutuamente capacidad jurídica y de obrar suficiente para
la celebración del presente Contrato.

\subsection{EXPONEN}\label{exponen}

PRIMERO: Que la CLIENTA tiene interés en la contratación de:

\begin{enumerate}
\def\labelenumi{\alph{enumi})}
\itemsep1pt\parskip0pt\parsep0pt
\item
  Servicio de mantenimiento de los sistemas ofrecidos e instalados (o
  revisados) por la PROVEEDORA en los exteriores de la vivienda para la
  conectividad con la red Guifi.net,
\item
  Acceso y uso de los servicios proporcionados desde los servidores de
  la PROVEEDORA,
\end{enumerate}

SEGUNDO: Que la PROVEEDORA es una entidad especializada en la prestación
de servicios de instalación, mantenimiento y conectividad sobre la red
Guifi.Net, cumpliendo con los requerimientos de la red abierta, libre y
neutral (Procomún RALN), aportando a la misma sus propios servicios en
línea.

TERCERO: Que las Partes están interesadas en celebrar un contrato de
Mantenimiento en virtud del cual la PROVEEDORA preste a la CLIENTA los
servicios de:

\begin{enumerate}
\def\labelenumi{\alph{enumi})}
\itemsep1pt\parskip0pt\parsep0pt
\item
  Mantenimiento de HARDWARE: Mantenimiento de Antenas y cableado,
  routers y switches, dispositivos de telefonía IP.
\item
  Mantenimiento de SOFTWARE: Mantenimiento de servicios ofrecidos en la
  red, y software de antenas y dispositivos de conectividad.
\item
  Mantenimiento de soporte de RED: Mantenimiento de la conectividad y la
  seguridad de la misma.
\end{enumerate}

Que las Partes acuerdan celebrar el presente contrato de MANTENIMIENTO Y
SERVICIOS, en adelante, el ``CONTRATO'', de acuerdo con las siguientes

\subsection{CLÁUSULAS}\label{cluxe1usulas}

\subsubsection{PRIMERA.- OBJETO}\label{primera.--objeto}

En virtud del Contrato la PROVEEDORA se obliga a prestar a la CLIENTA
los servicios de mantenimiento de hardware, software y de red
solicitados, en adelante los ``SERVICIOS'', en los términos y
condiciones previstos en el Contrato y en todos sus Anexos.

\subsubsection{SEGUNDA.- TÉRMINOS Y CONDICIONES GENERALES Y ESPECÍFICOS
DE PRESTACIÓN DE LOS
SERVICIOS}\label{segunda.--tuxe9rminos-y-condiciones-generales-y-especuxedficos-de-prestaciuxf3n-de-los-servicios}

2.1. Los Servicios se prestarán en los siguientes términos y condiciones
\textbf{generales}:

2.1.1. La PROVEEDORA responderá de la calidad del trabajo desarrollado
con la diligencia exigible a una empresa experta en la realización de
los trabajos objeto del Contrato.

2.1.2. La PROVEEDORA se obliga a gestionar y obtener, a su cargo, todas
las licencias, permisos y autorizaciones administrativas que pudieren
ser necesarias para la realización de los Servicios.

2.1.3. La PROVEEDORA se hará cargo de la totalidad de los tributos,
cualquiera que sea su naturaleza y carácter, que se devenguen como
consecuencia del Contrato, así como cualesquiera operaciones físicas y
jurídicas que conlleve, salvo el Impuesto sobre el Valor Añadido (IVA) o
su equivalente, que la PROVEEDORA repercutirá a la CLIENTA.

2.1.4. La PROVEEDORA guardará confidencialidad sobre la información que
le facilite la CLIENTA en o para la ejecución del Contrato o que por su
propia naturaleza deba ser tratada como tal. Se excluye de la categoría
de información confidencial toda aquella información que sea divulgada
por la CLIENTA, aquella que haya de ser revelada de acuerdo con las
leyes o con una resolución judicial o acto de autoridad competente,
durante el periodo máximo que estipule la ley.

2.1.5. En el caso de que la prestación de los Servicios suponga la
necesidad de acceder a datos de carácter personal, la PROVEEDORA, como
encargado del tratamiento, queda obligado al cumplimiento de la Ley
15/1999, de 13 de Diciembre, de Protección de Datos de Carácter Personal
y del Real Decreto 1720/2007, de 21 de diciembre, por el que se aprueba
el Reglamento de desarrollo de la Ley Orgánica 15/1999 y demás normativa
aplicable.

La PROVEEDORA responderá, por tanto, de las infracciones en que pudiera
incurrir en el caso de que destine los datos personales a otra
finalidad, los comunique a un tercero, o en general, los utilice de
forma irregular, así como cuando no adopte las medidas correspondientes
para el almacenamiento y custodia de los mismos. A tal efecto, se obliga
a indemnizar a la CLIENTA, por cualesquiera daños y perjuicios que sufra
directamente, o por toda reclamación, acción o procedimiento, que traiga
su causa de un incumplimiento o cumplimiento defectuoso por parte de la
PROVEEDORA de lo dispuesto tanto en el Contrato como lo dispuesto en la
normativa reguladora de la protección de datos de carácter personal.

A los efectos del artículo 12 de la Ley 15/1999, la PROVEEDORA
únicamente tratará los datos de carácter personal a los que tenga acceso
conforme a las instrucciones de la CLIENTA y no los aplicará o utilizará
con un fin distinto al objeto del Contrato, ni los comunicará, ni
siquiera para su conservación, a otras personas. En el caso de que la
PROVEEDORA destine los datos a otra finalidad, los comunique o los
utilice incumpliendo las estipulaciones del Contrato, será considerado
también responsable del tratamiento, respondiendo de las infracciones en
que hubiera incurrido personalmente.

La PROVEEDORA deberá adoptar las medidas de índole técnica y
organizativas necesarias que garanticen la seguridad de los datos de
carácter personal y eviten su alteración, pérdida, tratamiento o acceso
no autorizado, habida cuenta del estado de la tecnología, la naturaleza
de los datos almacenados y los riesgos a que están expuestos, ya
provengan de la acción humana o del medio físico o natural. A estos
efectos la PROVEEDORA deberá aplicar los niveles de seguridad que se
establecen en el Real Decreto 1720/2007 de acuerdo a la naturaleza de
los datos que trate.

2.1.6. La PROVEEDORA responderá de la corrección y precisión de los
documentos que aporte a la CLIENTA en ejecución del Contrato y avisará
sin dilación a la CLIENTA cuando detecte un error para que pueda adoptar
las medidas y acciones correctoras que estime oportunas.

2.1.7. Las obligaciones establecidas para la PROVEEDORA por la presente
cláusula serán también de obligado cumplimiento para sus posibles
empleados, colaboradores, tanto externos como internos, y
subcontratistas, por lo que la PROVEEDORA responderá frente a la CLIENTA
si tales obligaciones son incumplidas por tales empleados.

2.1.8. La CLIENTA podrá obtener de la PROVEEDORA las especificaciones de
su instalación, si lo requiere.

2.2. La PROVEEDORA prestará los Servicios en los siguientes términos y
condiciones \textbf{específicos}:

\begin{enumerate}
\def\labelenumi{\alph{enumi})}
\itemsep1pt\parskip0pt\parsep0pt
\item
  Mantenimiento de HARDWARE:

  \begin{itemize}
  \itemsep1pt\parskip0pt\parsep0pt
  \item
    \textbf{Mantenimiento preventivo} que incluye la monitorizacion
    remota detallada del correcto funcionamiento de los equipos de
    hardware instalados.
  \item
    \textbf{Mantenimiento correctivo} que incluye la sustitución de los
    componentes tanto internos como externos que fallaran en los
    equipos, todo ello a cargo de la CLIENTA, exceptuando los casos de
    defectos de fábrica o instalación defectuosa.
  \item
    \textbf{Mantenimiento evolutivo} que incluye la instalación de otros
    dispositivos adicionales para la mejora del rendimiento operativo en
    general y de seguridad. En el momento de la instalación, se
    determinarán la condiciones de propiedad y mantenimiento de estos
    nuevos elementos.
  \end{itemize}
\item
  Mantenimiento de SOFTWARE:

  \begin{itemize}
  \itemsep1pt\parskip0pt\parsep0pt
  \item
    \textbf{Mantenimiento preventivo} del software instalado que incluye
    la revisión de los parámetros críticos de los equipos y de la red.
  \item
    \textbf{Mantenimiento correctivo} que incluye la reinstalación o
    reconfiguración de software en el caso de anomalías en el
    funcionamiento.
  \item
    \textbf{Mantenimiento evolutivo} que incluye actualización e
    instalación de software adicional.
  \end{itemize}
\item
  Mantenimiento de RED:

  \begin{itemize}
  \itemsep1pt\parskip0pt\parsep0pt
  \item
    \textbf{Mantenimiento preventivo} que incluye la monitorización de
    los parámetros básicos de la red de forma que se garantice su
    adecuada dimensión.
  \item
    \textbf{Mantenimiento correctivo} que incluye la reconfiguración del
    hardware y software de RED después de la caída del sistema u otros
    percances.
  \item
    \textbf{Mantenimiento evolutivo} que incluye la configuración y
    optimización para el correcto funcionamiento de la comunicación
    entre los equipos conectados en RED, así como todos los periféricos
    conectados a ella que sean objeto de este contrato.
  \end{itemize}
\end{enumerate}

2.2.1. La PROVEEDORA se obliga a prestar los servicios necesarios para
el correcto funcionamiento del sistema instalado en la CLIENTA, siempre
que no se deban a manipulaciones indebidas. Dichos servicios comprenden
la realización de cuantas operaciones sean necesarias en los equipos o
sistemas incluidos en este contrato.

2.2.2. La CLIENTA facilitará la labor de la PROVEEDORA en todo momento
dándole el acceso a los sistemas requeridos. La PROVEEDORA se reserva el
acceso físico y telemático a los equipos comunes, a fin de poder
gestionar de forma eficiente y correcta tanto el apartado preventivo,
correctivo y/o evolutivo de estos sistemas.

2.2.3. El coste de cualquier componente que la PROVEEDORA tenga
sustituir o instalar será por cuenta de la CLIENTA, a excepción de
defectos de fábrica o mala manipulación por parte del PROVEEDOR o
terceros vinculados con el mismo. La PROVEEDORA no se responsabiliza de
los daños que pudieran producirse por la incorrecta manipulación de los
componentes realizados por la CLIENTA o terceros no vinculados.

2.2.4. La CLIENTA se compromete a utilizar los componentes del hardware
y del software de acuerdo con las instrucciones y el manual del
fabricante.

2.2.5. En ningún caso la PROVEEDORA mantendrá, instalará o configurará
sin los debidos permisos administrativos.

2.2.6. La PROVEEDORA ofrecerá un servicio de soporte técnico de 9 a 17
horas todos los días laborables no festivos mediante llamada telefónica
al (93-164-0492) o correo electrónico a soporte@guifibaix.coop. Las
vísperas de festivos nacionales, el horario de atención telefónica será
de 9 a 15 horas. Los fines de semana o festivos, el servicio estará
disponible a través del correo electrónico soporte@guifibaix.coop.
Cualquiera incidencia que se produzca, la CLIENTA deberá comunicarla a
la PROVEEDORA, a través del teléfono o correo electrónico mencionados,
para proceder a su solución.

\subsubsection{TERCERA.- POLÍTICA DE
USO}\label{tercera.--poluxedtica-de-uso}

3.1. La CLIENTA es el único responsable de determinar si los servicios
que constituyen el objeto de este Contrato se ajustan a sus necesidades,
por lo que la PROVEEDORA no garantiza que el servicio de mantenimiento
contratado se ajuste a las necesidades específicas de la CLIENTA.

3.2. La CLIENTA se obliga a hacer constar de forma clara, visible y
accesible desde sus contenidos, sus datos identificativos y como único
responsable de los contenidos, poniendo un aviso en sus contenidos de la
Política de uso de los mismos.

3.3. La CLIENTA se compromete a no hacer un uso de la red que altere por
sobrecarga o saturación, la calidad de servicio ofrecida al resto de
usuarios. La PROVEEDORA se reserva el derecho de bloquear o modular
dichos usos, para garantizar la calidad de servicio. La PROVEEDORA se
compromete a informar a la CLIENTA, si lo pide, si un servicio está
siendo modulado o bloqueado e informar de las alternativas.

3.4. La CLIENTA se compromete a no usar la red con fines delictivos, asi
como a no desarrollar actividades que persigan o tengan como
consecuencia:

\begin{itemize}
\itemsep1pt\parskip0pt\parsep0pt
\item
  La congestión de los enlaces de comunicaciones o sistemas
  informáticos.
\item
  La destrucción o modificación premeditada de la información de otros
  usuarios.
\item
  La violación de la privacidad e intimidad de otros usuarios.
\item
  El envío de SPAM a través de la red sea cual sea su forma
\end{itemize}

3.5. La PROVEEDORA no se hará responsable de las acciones o conductas de
la CLIENTA que usando la red o los servicios de la PROVEEDORA resultaran
ser delictivas o atentasen contra la Ley de Propiedad Intelectual o la
Ley Oficial de Protección de Datos.

3.6. Tanto CLIENTA como PROVEEDOR, se comprometen a subscribir las
condiciones de la Licencia de Red de ProComún, anexada a este contrato.

\subsubsection{CUARTA.- RESPONSABILIDAD
SOCIAL}\label{cuarta.--responsabilidad-social}

4.1. \textbf{Sin ánimo de lucro:} La PROVEEDORA se compromete a ajustar
sus precios y tarifas a la viabilidad económica del proyecto. Cualquier
superávit se reinvertirá en servicios o en infraestructuras comunes para
mejorar la red o se dedicará a fondos de reserva. Dentro de dicha
viabilidad económica, la PROVEEDORA tomará las decisiones bajo los
criterios sociales que se detallan a continuación.

4.2. \textbf{Red libre, abierta y neutral:} La PROVEEDORA se compromete
cumplir y hacer cumplir, en las redes que desplega y mantiene, la
licencia de Red Neutral y Abierta en Procomún anexada al presente
contrato. La PROVEEDORA se compromete a que los protocolos telemáticos
usados en dichas redes estén basados estándares abiertos.

4.3. \textbf{No atadura a la PROVEEDORA:} La PROVEEDORA se compromete a
que la CLIENTA pueda acceder, si lo pide, a la documentación técnica que
describe la arquitectura de la red, con el suficiente detalle para que
una persona o entidad especializada pueda substituir las funciones de la
PROVEEDORA. Los dispositivos instalados que consten en la factura como
pagados quedan en propiedad de la CLIENTA. La CLIENTA garantizará a la
PROVEEDORA el acceso físico y telemático a los dispositivos dentro de la
Red Libre y Neutral mientras estén conectados a ella y la PROVEEDORA lo
requiera para que la red siga funcionando, y a los dispositivos que
estén fuera de dicha red, mientras dure este contrato de mantenimiento,

4.4. \textbf{Software libre:} La PROVEEDORA se compromete con la CLIENTA
a que todos los aplicativos informáticos que desarrollen se publicarán
bajo una licencia libre según la definición de la Open Source
Initiative. También se compromete a desempeñar su actividad con relación
a la CLIENTA, usando exclusivamente software libre, exceptuando las
aplicaciones privativas para las que no hubiera alternativa libre
viable. En este caso, la PROVEEDORA se compromete a poner recursos para
viabilizar dicha alternativa.

4.5. \textbf{Formatos abiertos:} La PROVEEDORA se compromete a
comunicarse con la CLIENTA mediante documentos en formatos abiertos y
libres de regalías que no requieran la adquisición de licencias
privativas para acceder a ellos.

4.6. \textbf{Energías renovables:} La PROVEEDORA se compromete con la
CLIENTA a que los equipos y servidores localizados en sus sedes usen
exclusivamente energías renovables, contratándo la luz de dichas sedes
con comercializadoras eléctricas que garanticen el origen renovable de
la energía comprada.

4.7. \textbf{Reducción y tratamiento de residuos e impacto ambiental:}
La PROVEEDORA se compromete a reducir y reciclar o tratar adecuadamente
los residuos generados directamente de su labor. La PROVEEDORA se
compromete, si se lo pidiera la CLIENTA, a dar utilidad en la
instalación a los dispositivos electrónicos de los que ya disponga la
CLIENTA, si funcionalmente fuera viable. La PROVEEDORA facilitará la
recompra para su utilización en instalaciones nuevas, de los
dispositivos que, después de instalados, dejasen de usarse en la
instalación. La PROVEEDORA no esta obligado a dicha recompra ni a
hacerlo por el mismo precio que en la compra.

4.8. \textbf{Economía local:} La PROVEEDORA se compromete con la CLIENTA
a potenciar la economía local. La PROVEEDORA incluirá la localidad como
un criterio social más a tener en cuenta en la elección de sus
proveedores. La PROVEEDORA dará preferencia al pequeño comercio frente a
las grandes cadenas comerciales. La PROVEEDORA se compromete con la
CLIENTA a no deslocalizar trabajo remunerado, excluyendo explícitamente
el apoyo a proyectos de software libre o hardware libre, cuyos equipos
por su naturaleza son internacionales.

4.9. \textbf{Trabajo digno:} La PROVEEDORA se compromete a que los
trabajadores que dan el servicio a la CLIENTA tengan control democrático
sobre sus condiciones laborales. Los trabajadores serán remunerados en
igualdad a igual carga de trabajo, sin hacer discriminación por sexo,
orientación sexual, raza, discapacidad, origen social, religión o
ideología. La PROVEEDORA se compromete a mantener una horquilla de
sueldos por carga de trabajo inferior al 1:2, incluyendo en este
compromiso las subcontratas.

4.10. \textbf{Participación y formación:} La PROVEEDORA se compromete a
incorporar, mediante procesos democráticos, a los usuarios en las
decisiones que afecten a la evolución de la red común. Para tal fin, la
PROVEEDORA fomentará la autoorganizacion democrática de la CLIENTA con
otros clientes, y organizará acciones de formación técnica para que
estos puedan tomar decisiones informadas. La PROVEEDORA se compromete a
organizar acciones de formación de trabajadores y usuarios en el uso de
tecnologías libres.

4.11. \textbf{Extensión de los compromisos a terceros:} La PROVEEDORA
incidirá en terceros proveedores para que hagan extensivos los
compromisos sociales de la PROVEEDORA. Con tal fin, la PROVEEDORA se
compromete a dar preferencia a aquellos proveedores que cumplan dichos
compromisos. En casos en los que se tenga constancia o indicio
suficiente de que una tercera entidad proveedora no está respetando los
derechos humanos, los derechos de los trabajadores, los derechos de la
infancia o esté degradando el medio ambiente la PROVEEDORA se compromete
a, primero, intentar que cambie la situación, y si ello no se produce,
buscar terceros proveedores alternativos que sí los respeten.

\subsubsection{QUINTA.- PRECIO Y
FACTURACIÓN}\label{quinta.--precio-y-facturaciuxf3n}

Como contraprestación de los servicios prestados por la PROVEEDORA, la
CLIENTA abonará las tarifas adjuntas al presente contrato y que la
CLIENTA reconoce conocer y haber recibido en el momento de su firma.

La PROVEEDORA podrá variar las tarifas para adaptarlas por necesidades
del mercado, previa notificacion a la CLIENTA mínimo un mes antes de la
emisión de la primera factura en la que tengan efecto las nuevas
tarifas. Una vez notificadas las nuevas tarifas se entienden aceptadas
por la CLIENTA tácitamente.

La PROVEEDORA facturará, trimestralmente y a trimestre vista, a la
CLIENTA la cuota pactada incrementada con el IVA. Los pagos se llevarán
a cabo el dia 10 de los meses de enero, abril, julio y octubre o, en los
trimestres incompletos, el dia que comience a proporcionarse el
servicio. La PROVEEDORA facturará los periodos incompletos prorrateando
el número de días naturales efectivos respecto a los dias naturales del
trimestre.

En caso de impago de los cargos, la deuda vencida generará un interés de
demora igual al interés legal del dinero incrementado en dos puntos a
partir de la primera factura impagada y, si la factura no es pagada a
los dos meses, dará derecho a la PROVEEDORA a resolver el Contrato de
conformidad con lo dispuesto en las presentes Condiciones.

\subsubsection{SEXTA.- ACUERDO DE NIVEL DE
SERVICIO}\label{sexta.--acuerdo-de-nivel-de-servicio}

6.1. Todos los Servicios prestados por la PROVEEDORA se realizarán por
personal especializado en cada materia. El personal de la PROVEEDORA
atenderá remotamente o acudirá previsto de todo el material necesario,
adecuado y actualizado, para prestar los Servicios.

6.2. Las averías o el mal funcionamiento de los Servicios se comunicarán
a la PROVEEDORA a través de llamada telefónica (93-164-0492) o envío de
correo electrónico soporte@guifibaix.coop. Los horarios de atención, son
los especificados en el punto 2.2.6

6.3. Los problemas se atenderán en el período máximo establecido según
se defina la incidencia. El tiempo de resolución final, dependerá de la
complejidad de la misma, y del volumen de sistemas y/o usuarios
afectados.

6.4. Se entiende por \emph{incidencia crítica}: las incidencias que, en
el marco de la prestación de los Servicios, afectan significativamente a
la CLIENTA.

6.5. Se entiende por \emph{incidencia grave}: las incidencias que, en el
marco de la prestación de los Servicios, afectan moderadamente a la
CLIENTA.

6.6. Se entiende por \emph{incidencia leve}: las incidencias que se
limitan a entorpecer la prestación de los Servicios.

6.7. El tiempo de atención a dichas incidencias se realizará en los
siguientes períodos máximos desde el aviso:

\begin{itemize}
\itemsep1pt\parskip0pt\parsep0pt
\item
  Incidencia crítica: 30 minutos. (Tiempo máximo de resolución: 1 hora)
\item
  Incidencia grave: 1 hora.
\item
  Incidencia leve: entre 1 y 2 horas.
\end{itemize}

6.8. El estado de los Servicios se revisará semanalmente por la
PROVEEDORA para comprobar su correcto funcionamiento.

\subsubsection{SÉPTIMA.- DURACIÓN DEL
CONTRATO}\label{suxe9ptima.--duraciuxf3n-del-contrato}

El plazo de duración del presente Contrato es indefinido a partir de la
fecha referida en el encabezamiento del Contrato pudiendo el cliente a
proceder a su resolución en cualquier momento, efectuando un preaviso de
al menos un mes de antelación.

\subsubsection{OCTAVA.- SUSPENSION DEL
CONTRATO}\label{octava.--suspension-del-contrato}

En caso de que se hubiera producido un impago por parte de la CLIENTA, o
en caso de que la PROVEEDORA constate que la CLIENTA incumple alguna de
las clausulas de este Contrato, la PROVEEDORA, después de comunicarlo
por los canales especificados en este Contrato, podría suspender
temporalmente la prestación de los servicios que dependan de la
PROVEEDORA, excluyendo las garantías de conectividad con la Red Libre,
Abierta y Neutral.

La CLIENTA tiene derecho a solicitar y obtener de la PROVEEDORA la
suspensión temporal del servicio por un periodo determinado que no será
menor de un mes, ni superior a tres meses. El periodo de suspensión no
podrá exceder, en ningún caso, de noventa días por año natural. Durante
el tiempo que dure la suspensión temporal se facturará al cliente la
mitad del importe mensual, en concepto de mantenimiento de la red común.

\subsubsection{NOVENA.- MODIFICACIÓN DEL
CONTRATO}\label{novena.--modificaciuxf3n-del-contrato}

Las Partes podrán modificar el contrato de mutuo acuerdo y por escrito.

\subsubsection{DÉCIMA.- RESOLUCIÓN DEL
CONTRATO}\label{duxe9cima.--resoluciuxf3n-del-contrato}

Las Partes podrán resolver el Contrato, en caso de incumplimiento de las
obligaciones establecidas en el mismo.

\subsubsection{ONCEAVA.- NOTIFICACIONES}\label{onceava.--notificaciones}

Las notificaciones que se realicen las Partes deberán realizarse por
correo electrónico, teléfono o correo ordinario, a las siguientes
direcciones:

\begin{verbatim}
CLIENTA:
    Dirección:  C/Rue Percebe, 13, 4o 3a
                08970 Sant Joan Despí
    Email:      email@usuario.test
    Teléfono:   93-111-2222

PROVEEDOR:
    Dirección:  C/Major 22, 3
                08970 Sant Joan Despí
    Email:      contacto@guifibaix.coop
    Teléfono:   93-164-0492
\end{verbatim}

\subsubsection{DOCEAVA.- REGIMEN
JURÍDICO}\label{doceava.--regimen-juruxeddico}

El presente contrato tiene carácter mercantil, no existiendo en ningún
caso vínculo laboral alguno entre la CLIENTA y el personal de la
PROVEEDORA que preste concretamente los Servicios.

Toda controversia derivada de este contrato o que guarde relación con él
-- incluida cualquier cuestión relativa a su existencia, validez o
terminación - será resuelta mediante arbitraje DE DERECHO, administrado
por la \textbf{Asociación Europea de Arbitraje de Madrid (Aeade)}, de
conformidad con su Reglamento de Arbitraje vigente a la fecha de
presentación de la solicitud de arbitraje. El Tribunal Arbitral que se
designe a tal efecto estará compuesto por un único árbitro experto. La
sede del arbitraje será Barcelona.

~

Y en prueba de cuanto antecede, las Partes suscriben el Contrato, en dos
ejemplares y a un solo efecto, en el lugar y fecha señalados en el
encabezamiento.

~

\doubleSignature{ LA CLIENTA }{ LA PROVEEDORA }{ Sant Joan Despí }{20 de febrero de 2014}

\end{document}
